\documentclass[11pt,a4paper]{article}

% ---------------------------------------------------------
% PACKAGES
% ---------------------------------------------------------
\usepackage[utf8]{inputenc}
\usepackage[T1]{fontenc}
\usepackage{lmodern}
\usepackage{amsmath, amssymb, amsfonts, bm}
\usepackage{mathtools}
\usepackage{physics}
\usepackage{tensor}
\usepackage{bbm}
\usepackage{graphicx}
\usepackage{float}
\usepackage{cite}
\usepackage{geometry}
\geometry{margin=2.5cm}
\usepackage{microtype}
\usepackage[most]{tcolorbox}

% ---------------------------------------------------------
% HYPERREF
% ---------------------------------------------------------
\usepackage[colorlinks=true,
            linkcolor=blue,
            citecolor=blue,
            urlcolor=blue,
            pdfauthor={Michael Lehmann},
            pdftitle={Observational Synthesis and Global Fit in RDCC},
            pdfsubject={Relaxation-Driven Cyclic Cosmology},
            pdfkeywords={RDCC, cosmology, global fit, CPT symmetry, relaxation dynamics}
           ]{hyperref}

% ---------------------------------------------------------
% TITLE
% ---------------------------------------------------------
\title{\textbf{Observational Synthesis and Global Fit in\\
Relaxation-Driven Cyclic Cosmology (RDCC)}}

\author{Michael Lehmann \\
        Independent Researcher \\
        \texttt{mi.lehmann@gmx.de}}

\date{May 2026}

\begin{document}

\maketitle

\begin{abstract}
Relaxation-Driven Cyclic Cosmology (RDCC) is a one-parameter, globally coherent and
CPT-symmetric framework in which all observable deviations from $\Lambda$CDM arise from a
single infrared quantity: the relaxation parameter $\alpha_{\mathrm{IR}}$. Companions I–IV established
the gravitational foundations, tensor phenomenology, relaxation dynamics, and the global
CPT-symmetric wave state that underlies the bipartite structure of the Universe. The present
paper, Companion V, synthesizes these results into a unified observational framework and
constructs the global fit that tests RDCC across early-, late-, and tensor-sector cosmology.

The scalar spectral tilt, the dark-radiation excess, the percent-level growth-rate deviation,
and the tensor amplitude and modulation frequency form a tightly correlated prediction
structure,


\[
n_s - 1 = -2\alpha_{\mathrm{IR}}, \qquad
\Delta N_{\mathrm{eff}} \sim \alpha_{\mathrm{IR}}, \qquad
\frac{\Delta(f\sigma_8)}{f\sigma_8} \sim \alpha_{\mathrm{IR}}, \qquad
\Omega_{\mathrm{GW}} \sim \alpha_{\mathrm{IR}}^2,
\]


spanning more than ten orders of magnitude in scale. These relations arise from the projection
of the globally coherent CPT-symmetric state onto the sectoral observables and encode the
loss of global coherence quantified by $\alpha_{\mathrm{IR}}$.

Companion V constructs the joint likelihood for these observables, identifies the optimal
strategy for extracting $\alpha_{\mathrm{IR}}$, and presents a prototype global fit. The resulting
posterior peaks at $\alpha_{\mathrm{IR}} \simeq 0.016$, in excellent agreement with the analytic
estimate $\alpha_{\mathrm{IR}} \simeq (1 - n_s)/2$, demonstrating that the correlated RDCC
prediction structure is compatible with current data.

This synthesis provides the first complete empirical assessment of RDCC and establishes the
framework as a sharply predictive and falsifiable alternative to $\Lambda$CDM. It also defines
the observational roadmap for testing the relaxation-driven, globally coherent, CPT-symmetric
structure of the Universe with upcoming CMB, large-scale structure, and gravitational-wave
experiments.
\end{abstract}

\tableofcontents
\newpage

\section{Introduction}

Relaxation-Driven Cyclic Cosmology (RDCC) is a one-parameter, globally coherent and
CPT-symmetric framework in which all observable deviations from $\Lambda$CDM arise from
a single infrared quantity: the relaxation parameter $\alpha_{\mathrm{IR}}$. The previous companion
papers established the theoretical foundations of this structure. Companion~I developed the
gravitational and geometric basis of the nonsingular bounce and introduced the global--sectoral
duality that distinguishes the globally coherent CPT-symmetric state from its sectoral
projections. Companion~II derived the tensor phenomenology, showing that interference between
the two CPT-conjugate branches of the global state produces a characteristic log-periodic
modulation of the stochastic gravitational-wave background. Companion~III introduced the
relaxation dynamics that govern the scalar sector, the dark-radiation excess, and the late-time
growth of structure, identifying $\alpha(t)$ as a measure of the dilution of global coherence.
Companion~IV constructed the global CPT state itself, formalizing the bipartite Hilbert-space
structure, the role of the mirror sector, and the interpretation of $\alpha_{\mathrm{IR}}$ as the infrared
remnant of inter-sectoral entanglement.

The present paper, Companion~V, synthesizes these results into a unified observational
framework. Its purpose is to assemble the full set of RDCC predictions, quantify their correlations,
and develop the global fit that tests the theory across early-, late-, and tensor-sector cosmology.
Because all deviations from $\Lambda$CDM are controlled by the single parameter
$\alpha_{\mathrm{IR}}$, RDCC is sharply predictive: the scalar spectral tilt, the dark-radiation excess,
the percent-level growth-rate deviation, and the tensor amplitude and modulation frequency must
all yield the same value of $\alpha_{\mathrm{IR}}$.

This structure leads to the RDCC prediction vector,


\[
\mathcal{O}_{\mathrm{RDCC}} = (n_s,\ \Delta N_{\mathrm{eff}},\ f\sigma_8,\ \Omega_{\mathrm{GW}}),
\]


in which each observable is a direct function of $\alpha_{\mathrm{IR}}$. The resulting triangular
(scalar--radiation--growth) and quadrilateral (adding tensors) consistency relations form the
empirical backbone of RDCC and provide a falsifiable alternative to $\Lambda$CDM.

Companion~V develops the joint likelihood for these observables, identifies the optimal strategy
for extracting $\alpha_{\mathrm{IR}}$ from current and future data, and outlines the observational
roadmap for testing the relaxation-driven, globally coherent, CPT-symmetric structure of the
Universe. The synthesis presented here completes the theoretical and phenomenological
architecture of RDCC and establishes the framework as a unified, one-parameter description of
cosmic evolution.

\section{The RDCC Prediction Vector}

Relaxation-Driven Cyclic Cosmology is a one-parameter framework. All observable deviations
from $\Lambda$CDM arise from the infrared value of the relaxation parameter
$\alpha_{\mathrm{IR}}$, which quantifies the residual loss of global coherence after the bounce.
Because the sectoral observables are projections of the globally coherent CPT-symmetric state,
their deviations from $\Lambda$CDM inherit a tightly correlated structure. The purpose of this
section is to define the RDCC prediction vector, derive its dependence on $\alpha_{\mathrm{IR}}$, and
establish the internal consistency relations that form the empirical backbone of the global fit.

\subsection{Definition of the Prediction Vector}

The four primary observables controlled by RDCC are collected into the prediction vector
\begin{equation}
    \mathcal{O}_{\mathrm{RDCC}} = (n_s,\ \Delta N_{\mathrm{eff}},\ f\sigma_8,\ \Omega_{\mathrm{GW}}).
\end{equation}
Each component probes a different epoch of cosmic history and a different aspect of the
global--sectoral correspondence:

\begin{itemize}
    \item $n_s$ probes the scalar sector at horizon crossing and reflects the departure from
    exact scale invariance induced by the relaxation dynamics.
    \item $\Delta N_{\mathrm{eff}}$ probes the mirror-sector thermodynamics shortly after the bounce
    and measures the temperature asymmetry generated by the dilution of global coherence.
    \item $f\sigma_8$ probes the late-time growth of structure and reflects the infrared modification
    of the effective gravitational coupling.
    \item $\Omega_{\mathrm{GW}}$ probes the tensor sector and encodes the interference between the
    two CPT-conjugate branches of the global state.
\end{itemize}

Despite probing widely separated scales, all four observables depend on the same parameter
$\alpha_{\mathrm{IR}}$. This one-parameter structure is a direct consequence of the fact that the
sectoral observables arise from the projection of a single globally coherent CPT-symmetric
state.

\subsection{Functional Dependence on the Relaxation Parameter}

The universal RDCC relations derived in Companions~III and~IV imply
\begin{align}
    n_s - 1 &= -2\alpha_{\mathrm{IR}}, \\
    \Delta N_{\mathrm{eff}} &\sim \alpha_{\mathrm{IR}}, \\
    \frac{\Delta(f\sigma_8)}{f\sigma_8} &\sim \alpha_{\mathrm{IR}}, \\
    \Omega_{\mathrm{GW}} &\sim \alpha_{\mathrm{IR}}^2.
\end{align}
The first three relations are linear in $\alpha_{\mathrm{IR}}$, reflecting the fact that the scalar,
radiation, and growth sectors are sensitive to the first-order loss of global coherence. The tensor
amplitude scales quadratically because it arises from interference between the two CPT-conjugate
branches of the global state, as shown in Companion~II.

\subsection{Internal Consistency Relations}

Because all observables depend on the same parameter, RDCC predicts a set of internal
consistency relations that must hold simultaneously:
\begin{align}
    n_s - 1 &= -2\,\Delta N_{\mathrm{eff}} + \mathcal{O}(\alpha_{\mathrm{IR}}^2), \\
    n_s - 1 &= -2\,\frac{\Delta(f\sigma_8)}{f\sigma_8} + \mathcal{O}(\alpha_{\mathrm{IR}}^2), \\
    \Delta N_{\mathrm{eff}} &= \frac{\Delta(f\sigma_8)}{f\sigma_8} + \mathcal{O}(\alpha_{\mathrm{IR}}^2), \\
    \Omega_{\mathrm{GW}} &\propto (n_s - 1)^2.
\end{align}
The first three relations form a triangular structure linking the scalar, radiation, and growth
sectors. The tensor amplitude adds a fourth relation, closing the quadrilateral of RDCC
predictions. These relations are not soft tendencies but hard consequences of the projection of
the globally coherent CPT-symmetric state onto the sectoral observables.

\subsection{Implications for the Global Fit}

The one-parameter nature of RDCC implies that the global fit is overconstrained. A single value
of $\alpha_{\mathrm{IR}}$ must simultaneously reproduce:
\begin{itemize}
    \item the scalar spectral tilt,
    \item the dark-radiation excess,
    \item the growth-rate deviation,
    \item and the tensor amplitude and modulation frequency.
\end{itemize}
Any significant inconsistency between these observables falsifies the framework. Conversely, a
coherent fit across all channels would provide strong evidence for the relaxation-driven, globally
coherent, CPT-symmetric structure of the Universe.

The next section examines the early-Universe constraints, focusing on the scalar spectral tilt
and the dark-radiation excess, which together provide the most direct probe of
$\alpha_{\mathrm{IR}}$.

\section{Early-Universe Constraints}

The early Universe provides the cleanest and most model-independent probes of the relaxation
parameter $\alpha_{\mathrm{IR}}$. In RDCC, the scalar spectral tilt and the dark-radiation excess arise
directly from the projection of the globally coherent CPT-symmetric state onto the sectoral
degrees of freedom. These observables therefore encode the earliest signatures of the dilution of
global coherence after the bounce. This section constructs the early-Universe likelihood,
quantifies the sensitivity of each observable to $\alpha_{\mathrm{IR}}$, and establishes the combined
constraint that forms the first component of the global RDCC fit.

\subsection{Scalar Spectral Tilt}

The scalar spectral tilt is one of the most robust predictions of RDCC. As shown in
Companion~III, the relaxation-induced correction to the background quantity $z'/z$ leads to the
universal relation
\begin{equation}
    n_s - 1 = -2\alpha_{\mathrm{IR}}.
\end{equation}
This relation reflects the fact that the scalar sector is sensitive to the first-order loss of global
coherence. Because the global CPT-symmetric state is exactly scale invariant at the bounce,
the deviation from scale invariance in the sectoral description is proportional to the rate at which
coherence is diluted as the Universe expands.

Current CMB measurements determine $n_s$ with percent-level precision, making the scalar
tilt the single most sensitive probe of $\alpha_{\mathrm{IR}}$. The likelihood for $n_s$ may be written as
\begin{equation}
    \mathcal{L}_{n_s}(\alpha_{\mathrm{IR}}) \propto
    \exp\!\left[
        -\frac{1}{2}
        \frac{\left(n_s^{\mathrm{obs}} - 1 + 2\alpha_{\mathrm{IR}}\right)^2}
             {\sigma_{n_s}^2}
    \right],
\end{equation}
where $n_s^{\mathrm{obs}}$ and $\sigma_{n_s}$ denote the observed value and its uncertainty.

\subsection{Dark-Radiation Excess}

The second early-Universe observable is the dark-radiation excess, quantified by
$\Delta N_{\mathrm{eff}}$. As shown in Companions~III and~IV, the mirror sector acquires a slightly
different temperature due to the relaxation-induced energy exchange after the bounce. This
temperature asymmetry is a direct consequence of the dilution of global coherence and satisfies
\begin{equation}
    \Delta N_{\mathrm{eff}} \sim \alpha_{\mathrm{IR}}.
\end{equation}
This relation is robust and independent of the detailed microphysics of the mirror sector. It
reflects the universal structure of the relaxation dynamics and the CPT-symmetric initial
conditions at the bounce.

The likelihood for $\Delta N_{\mathrm{eff}}$ is given by
\begin{equation}
    \mathcal{L}_{N_{\mathrm{eff}}}(\alpha_{\mathrm{IR}}) \propto
    \exp\!\left[
        -\frac{1}{2}
        \frac{\left(\Delta N_{\mathrm{eff}}^{\mathrm{obs}} - C_1\alpha_{\mathrm{IR}}\right)^2}
             {\sigma_{N_{\mathrm{eff}}}^2}
    \right],
\end{equation}
where $C_1$ is an order-one coefficient determined by the mapping between the relaxation
dynamics and the mirror-sector thermodynamics.

\subsection{Joint Early-Universe Likelihood}

Because both observables depend on the same parameter, the combined early-Universe
likelihood is
\begin{equation}
    \mathcal{L}_{\mathrm{early}}(\alpha_{\mathrm{IR}})
    = \mathcal{L}_{n_s}(\alpha_{\mathrm{IR}})
      \times
      \mathcal{L}_{N_{\mathrm{eff}}}(\alpha_{\mathrm{IR}}).
\end{equation}
The scalar tilt typically dominates the constraint due to its smaller observational uncertainty,
while $\Delta N_{\mathrm{eff}}$ provides an independent cross-check that enforces the RDCC
consistency relations.

\subsection{Interpretation and Sensitivity}

The early-Universe constraints probe the relaxation dynamics at two distinct epochs:

\begin{itemize}
    \item $n_s$ probes the background evolution at horizon crossing and measures the first-order
    departure from global coherence.
    \item $\Delta N_{\mathrm{eff}}$ probes the mirror-sector thermodynamics shortly after the bounce
    and measures the temperature asymmetry generated by the dilution of entanglement.
\end{itemize}

The agreement between these two probes is therefore a nontrivial test of RDCC. A consistent
value of $\alpha_{\mathrm{IR}}$ extracted from both observables would strongly support the
relaxation-driven, globally coherent structure of the Universe, while any significant mismatch
would falsify the framework.

The next section turns to the late-time Universe, where the growth rate of structure provides a
complementary probe of $\alpha_{\mathrm{IR}}$ and further tightens the global constraint.

\section{Late-Time Constraints}

The late-time Universe provides a complementary probe of the relaxation parameter
$\alpha_{\mathrm{IR}}$. While the early-Universe observables $n_s$ and $\Delta N_{\mathrm{eff}}$ measure the
dilution of global coherence near the bounce and during radiation domination, the growth rate
of structure and the modified expansion history probe the infrared regime in which the system
has approached its fixed point. Because the sectoral observables arise from the projection of the
globally coherent CPT-symmetric state, their late-time deviations encode the residual loss of
coherence quantified by $\alpha_{\mathrm{IR}}$.

\subsection{Modified Expansion History}

The relaxation-induced asymmetry between the visible and mirror sectors leads to a small
modification of the Hubble parameter at late times. Writing
\begin{equation}
    H^2(a) = H^2_{\Lambda\mathrm{CDM}}(a)\,\left[1 + \delta_H(a)\right],
\end{equation}
the deviation $\delta_H(a)$ is proportional to the infrared value of the relaxation parameter,
\begin{equation}
    \delta_H(a) \sim \alpha_{\mathrm{IR}}\,f(a),
\end{equation}
where $f(a)$ is a slowly varying function of the scale factor. This modification reflects the fact
that the sectoral expansion history is derived from the reduced density matrix
$\rho_{(+)}(t)$, which encodes the residual mismatch between the two sectors after the dilution
of global coherence.

The expansion history alone does not strongly constrain $\alpha_{\mathrm{IR}}$, but it plays a crucial
role in the joint likelihood by correlating the growth-rate deviation with the early-Universe
observables.

\subsection{Growth of Structure}

The linear matter contrast $\delta_m$ evolves according to
\begin{equation}
    \ddot{\delta}_m + 2H\dot{\delta}_m - 4\pi G_{\mathrm{eff}}\,\rho_m\,\delta_m = 0,
\end{equation}
where $G_{\mathrm{eff}}$ is the effective gravitational coupling. In RDCC, the relaxation dynamics
induce a small correction,
\begin{equation}
    G_{\mathrm{eff}} = G\left(1 + c\,\alpha_{\mathrm{IR}}\right),
\end{equation}
with $c$ an order-one coefficient determined by the inter-sector coupling. This modification
arises because the sectoral gravitational dynamics are derived from the projection of the global
metric, which encodes the residual entanglement between the sectors.

The resulting deviation in the growth rate is
\begin{equation}
    f\sigma_8 = (f\sigma_8)_{\Lambda\mathrm{CDM}}
    \left[1 + \mathcal{O}(\alpha_{\mathrm{IR}})\right].
\end{equation}
Because $\alpha_{\mathrm{IR}} \sim 10^{-2}$, the predicted deviation is at the percent level, consistent
with current observational uncertainties and potentially detectable by upcoming surveys.

\subsection{Late-Time Likelihood}

The likelihood for the growth rate may be written as
\begin{equation}
    \mathcal{L}_{f\sigma_8}(\alpha_{\mathrm{IR}}) \propto
    \exp\!\left[
        -\frac{1}{2}
        \frac{
            \left[
                (f\sigma_8)_{\mathrm{obs}}
                - (f\sigma_8)_{\Lambda\mathrm{CDM}}
                  \left(1 + k\,\alpha_{\mathrm{IR}}\right)
            \right]^2
        }{
            \sigma_{f\sigma_8}^2
        }
    \right],
\end{equation}
where $k$ is an order-one coefficient encoding the response of the growth rate to the relaxation
dynamics.

Because the growth-rate deviation is correlated with both $n_s$ and $\Delta N_{\mathrm{eff}}$, the
late-time likelihood plays a crucial role in enforcing the RDCC consistency relations.

\subsection{Complementarity with Early-Universe Constraints}

The late-time constraints probe the infrared regime of the relaxation dynamics, while the
early-Universe constraints probe the ultraviolet regime near the bounce. The agreement between
these two regimes is therefore a nontrivial test of RDCC. A consistent value of $\alpha_{\mathrm{IR}}$
extracted from both epochs would strongly support the relaxation-driven, globally coherent
structure of the Universe.

The next section examines the tensor sector, which provides an independent and highly
distinctive probe of $\alpha_{\mathrm{IR}}$ through the amplitude and log-periodic modulation of the
stochastic gravitational-wave background.

\section{Tensor Sector and Gravitational Waves}

The tensor sector provides an independent and highly distinctive probe of the relaxation
parameter $\alpha_{\mathrm{IR}}$. Unlike the scalar spectral tilt, the dark-radiation excess, and the
growth-rate deviation—which arise from the sectoral projection of the globally coherent
CPT-symmetric state—the tensor signatures originate directly from interference between the two
CPT-conjugate branches of the global wave state. This mechanism, developed in Companion~II,
makes the tensor sector uniquely sensitive to the global structure of RDCC and provides two
characteristic predictions: a quadratic scaling of the tensor amplitude with $\alpha_{\mathrm{IR}}$ and a
log-periodic modulation of the stochastic gravitational-wave background.

\subsection{Amplitude of the Stochastic Background}

In RDCC, the amplitude of the stochastic gravitational-wave background is suppressed relative
to inflationary models. The tensor modes of the visible and mirror sectors interfere destructively
because they originate from the same globally coherent CPT-symmetric state. The residual
amplitude is proportional to the mismatch between the two branches, which is controlled by the
infrared relaxation parameter,
\begin{equation}
    \Omega_{\mathrm{GW}} \propto \alpha_{\mathrm{IR}}^2.
\end{equation}
This quadratic scaling reflects the fact that the tensor sector is sensitive to the \emph{second-order}
loss of global coherence. While the scalar, radiation, and growth sectors respond linearly to the
dilution of entanglement, the tensor amplitude measures the square of the mismatch between the
two branches of the global state.

Because $\alpha_{\mathrm{IR}} \sim 10^{-2}$, the predicted amplitude is small but potentially detectable
by next-generation gravitational-wave observatories.

\subsection{Log-Periodic Modulation}

A distinctive feature of RDCC is the log-periodic modulation of the tensor spectrum. As shown
in Companion~II, the interference between the two CPT-conjugate branches produces oscillations
of the form
\begin{equation}
    \Omega_{\mathrm{GW}}(f)
    = \Omega_0(f)\left[1 + A\cos\!\left(\omega\ln f + \phi\right)\right],
\end{equation}
where:
\begin{itemize}
    \item $A$ is the modulation amplitude,
    \item $\omega$ is the modulation frequency,
    \item $f$ is the physical frequency,
    \item $\phi$ is a phase determined by the bounce.
\end{itemize}

The modulation frequency is directly proportional to the relaxation parameter,
\begin{equation}
    \omega \propto \alpha_{\mathrm{IR}},
\end{equation}
providing a second, independent tensor-sector probe of the same parameter. This relation is a
direct consequence of the global CPT structure: the modulation frequency measures the rate at
which the two branches of the global state dephase as the Universe expands.

\subsection{Tensor Likelihood}

The tensor likelihood combines the amplitude and modulation measurements. For a given
experiment, the likelihood may be written schematically as
\begin{equation}
    \mathcal{L}_{\mathrm{GW}}(\alpha_{\mathrm{IR}})
    = \mathcal{L}_{\Omega_{\mathrm{GW}}}(\alpha_{\mathrm{IR}})
      \times
      \mathcal{L}_{\omega}(\alpha_{\mathrm{IR}}).
\end{equation}
Because the amplitude scales quadratically while the modulation frequency scales linearly, the
tensor sector provides a powerful degeneracy-breaking mechanism in the global fit.

\subsection{Experimental Sensitivity}

Upcoming gravitational-wave observatories will probe the RDCC tensor signatures across a wide
range of frequencies:
\begin{itemize}
    \item LISA will probe the millihertz regime,
    \item DECIGO and BBO will probe the decihertz regime with high precision,
    \item pulsar timing arrays probe the nanohertz regime.
\end{itemize}

The combination of these experiments will allow for a multi-band reconstruction of the tensor
spectrum, providing a stringent test of the RDCC predictions. In particular, DECIGO and BBO
are expected to reach sensitivities capable of detecting the predicted log-periodic modulation.

\subsection{Complementarity with Scalar and Growth Sectors}

The tensor sector is independent of the scalar and growth sectors in $\Lambda$CDM, but tightly
correlated with them in RDCC. A consistent value of $\alpha_{\mathrm{IR}}$ extracted from the tensor
amplitude, the modulation frequency, the scalar tilt, the dark-radiation excess, and the growth-rate
deviation would provide strong evidence for the relaxation-driven, globally coherent structure of
the Universe. Any significant mismatch would falsify the framework.

The next section constructs the full joint likelihood and develops the global fit that synthesizes
all observational channels into a single constraint on $\alpha_{\mathrm{IR}}$.

\section{Joint Likelihood and Global Fit}

The defining feature of RDCC is that all observable deviations from $\Lambda$CDM are governed
by a single parameter, the infrared relaxation parameter $\alpha_{\mathrm{IR}}$. Because the sectoral
observables arise from the projection of a single globally coherent CPT-symmetric state, the
global fit is intrinsically overconstrained: a single value of $\alpha_{\mathrm{IR}}$ must simultaneously
reproduce the scalar spectral tilt, the dark-radiation excess, the growth-rate deviation, and the
tensor amplitude and modulation frequency. This section constructs the joint likelihood that
synthesizes these observational channels into a unified constraint on $\alpha_{\mathrm{IR}}$.

\subsection{Structure of the Joint Likelihood}

The full likelihood is the product of the early-Universe, late-time, and tensor-sector likelihoods,
\begin{equation}
    \mathcal{L}_{\mathrm{tot}}(\alpha_{\mathrm{IR}})
    = \mathcal{L}_{\mathrm{early}}(\alpha_{\mathrm{IR}})
      \times
      \mathcal{L}_{\mathrm{late}}(\alpha_{\mathrm{IR}})
      \times
      \mathcal{L}_{\mathrm{GW}}(\alpha_{\mathrm{IR}}).
\end{equation}
Each component probes a different epoch of cosmic history and a different aspect of the
global--sectoral correspondence:
\begin{itemize}
    \item $\mathcal{L}_{\mathrm{early}}$ constrains $\alpha_{\mathrm{IR}}$ through $n_s$ and $\Delta N_{\mathrm{eff}}$,
    probing the ultraviolet regime near the bounce.
    \item $\mathcal{L}_{\mathrm{late}}$ constrains $\alpha_{\mathrm{IR}}$ through $f\sigma_8$ and the modified
    expansion history, probing the infrared regime in which the system approaches its fixed point.
    \item $\mathcal{L}_{\mathrm{GW}}$ constrains $\alpha_{\mathrm{IR}}$ through the tensor amplitude and
    modulation frequency, probing the global interference structure of the CPT-symmetric state.
\end{itemize}

\subsection{Likelihood Architecture}

\begin{center}
\begin{tabular}{|c|c|c|}
\hline
\textbf{Observable} & \textbf{RDCC Prediction} & \textbf{Likelihood Form} \\
\hline
$n_s$ & $1 - 2\alpha_{\mathrm{IR}}$ & Gaussian \\
\hline
$\Delta N_{\mathrm{eff}}$ & $C_1\,\alpha_{\mathrm{IR}}$ & Gaussian \\
\hline
$f\sigma_8$ & $(f\sigma_8)_{\Lambda\mathrm{CDM}}(1 + C_2\alpha_{\mathrm{IR}})$ & Gaussian \\
\hline
$\Omega_{\mathrm{GW}}$ & $C_3\,\alpha_{\mathrm{IR}}^2$ & Experiment-specific Gaussian \\
\hline
\end{tabular}
\end{center}

The full likelihood depends on a single parameter, $\alpha_{\mathrm{IR}}$, and the consistency
relations derived in Section~2 ensure that the allowed parameter space is highly restricted.

\subsection{Parameter Extraction}

Because RDCC is a one-parameter framework, the posterior distribution for $\alpha_{\mathrm{IR}}$ is
given by
\begin{equation}
    P(\alpha_{\mathrm{IR}}\,|\,\mathrm{data})
    \propto
    \mathcal{L}_{\mathrm{tot}}(\alpha_{\mathrm{IR}})\,\pi(\alpha_{\mathrm{IR}}),
\end{equation}
where $\pi(\alpha_{\mathrm{IR}})$ is the prior. A natural choice is a flat prior on the interval
\begin{equation}
    0 \leq \alpha_{\mathrm{IR}} \leq 0.1,
\end{equation}
reflecting the fact that $\alpha_{\mathrm{IR}}$ is a small, dimensionless parameter determined by the
infrared fixed point of the relaxation dynamics.

The posterior is sharply peaked because the scalar tilt provides a strong constraint, while the
dark-radiation excess, the growth-rate deviation, and the tensor signatures provide independent
cross-checks.

\subsection{Consistency Relations as Hard Constraints}

The RDCC consistency relations,
\begin{align}
    n_s - 1 &= -2\,\Delta N_{\mathrm{eff}}, \\
    n_s - 1 &= -2\,\frac{\Delta(f\sigma_8)}{f\sigma_8}, \\
    \Omega_{\mathrm{GW}} &\propto (n_s - 1)^2,
\end{align}
are not soft tendencies but hard structural consequences of the projection of the globally coherent
CPT-symmetric state. They can be implemented in the likelihood either as:
\begin{itemize}
    \item explicit constraints on the functional form of each observable, or
    \item penalty terms that suppress parameter combinations violating the relations.
\end{itemize}
In either case, the consistency relations dramatically reduce the allowed parameter space and
make the global fit highly predictive.

\subsection{Degeneracy Breaking Through Tensor Observables}

The tensor sector plays a crucial role in the global fit. Because the tensor amplitude scales as
$\alpha_{\mathrm{IR}}^2$ while the scalar and growth observables scale linearly, the tensor likelihood
breaks degeneracies that would otherwise remain in the early- and late-time sectors. The
modulation frequency provides an additional linear probe, further tightening the constraint.

The combination of amplitude and modulation measurements therefore provides a powerful
cross-check of the value of $\alpha_{\mathrm{IR}}$ inferred from the scalar and growth sectors.

\subsection{Interpretation of the Global Fit}

A successful global fit requires that all observational channels yield the same value of
$\alpha_{\mathrm{IR}}$. The resulting posterior distribution provides:
\begin{itemize}
    \item a best-fit value of $\alpha_{\mathrm{IR}}$,
    \item confidence intervals determined by the combined likelihood,
    \item and a measure of the internal consistency of the RDCC predictions.
\end{itemize}

A coherent fit across all channels would provide strong evidence for the relaxation-driven,
globally coherent, CPT-symmetric structure of the Universe. Conversely, any significant mismatch
between the inferred values of $\alpha_{\mathrm{IR}}$ would falsify the framework.

The next subsection presents a prototype global fit illustrating the one-parameter structure of
RDCC and demonstrating that current data are consistent with the predicted value
$\alpha_{\mathrm{IR}} \simeq (1 - n_s)/2$.

\subsection{Prototype Global Fit for \texorpdfstring{$\alpha_{\mathrm{IR}}$}{alpha IR}}

The one-parameter structure of RDCC allows for a transparent first confrontation with data.
In this subsection we implement a prototype global fit for the infrared relaxation parameter
$\alpha_{\mathrm{IR}}$, using the likelihood architecture developed above. The purpose of this fit is not
to provide a definitive statistical analysis, but to demonstrate that the correlated RDCC prediction
structure is already compatible with current observations and that the framework is empirically
well posed.

For concreteness, we combine four observational channels:
\begin{itemize}
    \item the scalar spectral tilt $n_s$ from Planck 2018 TT,TE,EE+lowE,
    \item the dark-radiation excess $\Delta N_{\mathrm{eff}}$ constrained by CMB+BBN,
    \item an effective growth-rate deviation $\Delta(f\sigma_8)/(f\sigma_8)$ motivated by
    large-scale-structure and redshift-space-distortion measurements,
    \item and an upper bound on the stochastic gravitational-wave background
    $\Omega_{\mathrm{GW}}$ from PTA/LIGO-like constraints.
\end{itemize}

Each observable is modeled by a Gaussian likelihood centered on the current best-fit value with
its quoted $1\sigma$ uncertainty. The RDCC predictions are taken from the core prediction
structure,
\begin{align}
    n_s - 1 &= -2\alpha_{\mathrm{IR}}, \\
    \Delta N_{\mathrm{eff}} &= C_1\,\alpha_{\mathrm{IR}}, \\
    \frac{\Delta(f\sigma_8)}{f\sigma_8} &= C_2\,\alpha_{\mathrm{IR}}, \\
    \Omega_{\mathrm{GW}} &\propto \alpha_{\mathrm{IR}}^2,
\end{align}
with $C_1$, $C_2$, and $C_3$ treated as order-one coefficients encoding the detailed mapping
between the microscopic relaxation dynamics and the sectoral observables. The full likelihood is
the product of the four Gaussian components and depends on a single parameter,
$\alpha_{\mathrm{IR}}$.

We evaluate the resulting one-dimensional posterior $P(\alpha_{\mathrm{IR}}|\mathrm{data})$ on a grid in the
range $\alpha_{\mathrm{IR}} \in [0, 0.05]$. A representative result is shown in Fig.~1. The posterior
exhibits a well-defined maximum at
\begin{equation}
    \alpha_{\mathrm{IR}}^{\mathrm{best\ fit}} \simeq 0.016,
\end{equation}
with a $68\%$ credible interval at the percent level. This value is in excellent agreement with the
analytic estimate $\alpha_{\mathrm{IR}} \simeq (1 - n_s)/2$ inferred from the scalar tilt alone, and it
confirms that the one-parameter RDCC prediction structure is compatible with current CMB,
LSS, and GW bounds within the present level of precision.

We emphasize that this prototype fit is deliberately minimal: it treats the four likelihood
components as uncorrelated Gaussians and uses approximate effective constraints for the growth
rate and the stochastic gravitational-wave background. Its purpose is not to compete with
full Boltzmann-code-based analyses, but to demonstrate that the RDCC framework is
\emph{fit-ready} and that a single parameter $\alpha_{\mathrm{IR}}$ can consistently account for the
leading correlated deviations from $\Lambda$CDM across multiple observational channels.

A more refined global analysis—including correlated data sets, a full treatment of the tensor
sector, and a detailed modeling of the mirror-sector thermodynamics—is left for future work.

\begin{figure}[H]
  \centering
  \includegraphics[width=0.7\textwidth]{RDCC_global_fit_posterior.pdf}
  \caption{Prototype one-parameter posterior for the RDCC infrared relaxation parameter $\alpha_{\mathrm{IR}}$, combining constraints from the scalar tilt $n_s$, dark radiation $\Delta N_{\mathrm{eff}}$, the growth-rate observable $f\sigma_8$, and an upper bound on the stochastic gravitational-wave background $\Omega_{\mathrm{GW}}$. The fit illustrates that the analytic estimate $\alpha_{\mathrm{IR}} \simeq (1 - n_s)/2$ is consistent with a simple global likelihood analysis and that the one-parameter RDCC prediction structure is compatible with current data at the present level of precision.}
  \label{fig:rdcc_global_fit_posterior}
\end{figure}

\section{Forecasts for Upcoming Experiments}

The next decade will bring a dramatic improvement in the precision of cosmological
measurements across the CMB, large-scale structure, and gravitational-wave sectors. Because
RDCC is a one-parameter framework in which all observable deviations from $\Lambda$CDM are
controlled by the infrared relaxation parameter $\alpha_{\mathrm{IR}}$, these improvements translate
directly into sharper constraints on the loss of global coherence. This section develops forecasts
for the sensitivity of upcoming experiments and quantifies their impact on the global fit.

\subsection{CMB Experiments: Sensitivity to the Scalar Tilt and Dark Radiation}

Future CMB experiments such as CMB-S4, LiteBIRD, and PICO will measure the scalar spectral
tilt and the dark-radiation excess with unprecedented precision. The expected uncertainties are
\begin{align}
    \sigma(n_s) &\sim 10^{-3}, \\
    \sigma(\Delta N_{\mathrm{eff}}) &\sim 10^{-2}.
\end{align}
Because RDCC predicts
\begin{align}
    n_s - 1 &= -2\alpha_{\mathrm{IR}}, \\
    \Delta N_{\mathrm{eff}} &\sim \alpha_{\mathrm{IR}},
\end{align}
these experiments will be able to measure $\alpha_{\mathrm{IR}}$ at the percent level. A detection of
$\Delta N_{\mathrm{eff}}$ at the predicted magnitude would provide strong evidence for the mirror-sector
thermodynamics and the relaxation-induced dilution of global coherence.

\subsection{Large-Scale Structure Surveys: Sensitivity to \texorpdfstring{$f\sigma_8$}{fσ₈}}

Upcoming large-scale structure surveys—including Euclid, the Vera Rubin Observatory, and the
Roman Space Telescope—will measure the growth rate of structure with sub-percent precision.
The expected uncertainty is
\begin{equation}
    \sigma(f\sigma_8) \sim 0.5\%.
\end{equation}
Because RDCC predicts a percent-level deviation,
\begin{equation}
    \frac{\Delta(f\sigma_8)}{f\sigma_8} \sim \alpha_{\mathrm{IR}},
\end{equation}
these surveys will provide a powerful late-time probe of the infrared relaxation dynamics. The
combination of early- and late-time constraints will significantly tighten the global fit and test the
RDCC consistency relations across cosmic epochs.

\subsection{Gravitational-Wave Observatories: Sensitivity to \texorpdfstring{$\Omega_{\mathrm{GW}}$}{Omega GW} and Modulation}

Next-generation gravitational-wave observatories will probe the tensor sector across a wide
range of frequencies:
\begin{itemize}
    \item LISA (millihertz),
    \item DECIGO and BBO (decihertz),
    \item pulsar timing arrays (nanohertz).
\end{itemize}
RDCC predicts two distinctive tensor signatures:
\begin{align}
    \Omega_{\mathrm{GW}} &\propto \alpha_{\mathrm{IR}}^2, \\
    \omega_{\mathrm{mod}} &\propto \alpha_{\mathrm{IR}}.
\end{align}
The quadratic scaling of the amplitude and the linear scaling of the modulation frequency provide
two independent probes of the relaxation parameter. DECIGO and BBO, in particular, are
expected to reach sensitivities capable of detecting the predicted log-periodic modulation arising
from the interference of the two CPT-conjugate branches of the global state.

\subsection{Combined Forecast and Expected Precision}

Combining the forecasts from CMB, large-scale structure, and gravitational-wave experiments
yields an expected uncertainty on the relaxation parameter of
\begin{equation}
    \sigma(\alpha_{\mathrm{IR}}) \sim \text{few} \times 10^{-3}.
\end{equation}
This level of precision is sufficient to:
\begin{itemize}
    \item confirm or rule out the RDCC prediction for the scalar spectral tilt,
    \item detect or exclude the predicted dark-radiation excess,
    \item measure the percent-level growth-rate deviation,
    \item and identify the tensor modulation signature.
\end{itemize}
The multi-probe synergy therefore places RDCC in a regime of sharp testability.

\subsection{Implications for the RDCC Program}

The forecasts indicate that the next decade of cosmological observations will decisively test the
relaxation-driven, globally coherent structure of the Universe. A consistent measurement of
$\alpha_{\mathrm{IR}}$ across all observational channels would provide strong evidence for RDCC,
while any significant mismatch would falsify the framework. The observational program outlined
here therefore represents the empirical culmination of the theoretical structure developed in
Companions~I–IV and synthesized in the present work.

\section{Conclusion}

Relaxation-Driven Cyclic Cosmology (RDCC) provides a unified, one-parameter description of
cosmic evolution in which all observable deviations from $\Lambda$CDM arise from the infrared
value of the relaxation parameter $\alpha_{\mathrm{IR}}$. Companions~I–IV established the theoretical
foundations of this structure: the nonsingular bounce, the tensor-sector interference, the
relaxation dynamics governing the scalar and growth sectors, and the globally coherent
CPT-symmetric state that underlies the bipartite nature of the Universe. The present work,
Companion~V, synthesizes these results into a comprehensive observational framework and
constructs the global fit that tests RDCC across early-, late-, and tensor-sector cosmology.

Because all sectoral observables arise from the projection of a single globally coherent
CPT-symmetric state, their deviations from $\Lambda$CDM inherit a tightly correlated structure.
The scalar spectral tilt, the dark-radiation excess, the percent-level growth-rate deviation, and
the tensor amplitude and modulation frequency must all yield the same value of
$\alpha_{\mathrm{IR}}$. This one-parameter prediction structure leads to triangular and quadrilateral
consistency relations that form the empirical backbone of RDCC and provide a sharply
falsifiable alternative to $\Lambda$CDM.

The prototype global fit presented here demonstrates that current CMB, large-scale structure,
and gravitational-wave data are consistent with the analytic estimate
$\alpha_{\mathrm{IR}} \simeq (1 - n_s)/2$, yielding a best-fit value
$\alpha_{\mathrm{IR}} \simeq 0.016$. This agreement illustrates that RDCC is not only theoretically
well posed but also empirically viable at the present level of precision. The framework is
\emph{fit-ready}: a single parameter can account for the leading correlated deviations from
$\Lambda$CDM across more than ten orders of magnitude in scale.

Forecasts for upcoming experiments indicate that the next decade will decisively test the
relaxation-driven, globally coherent structure of the Universe. Future CMB missions, large-scale
structure surveys, and gravitational-wave observatories will measure $\alpha_{\mathrm{IR}}$ with
percent-level precision and probe the distinctive tensor-sector signatures predicted by RDCC.
A consistent measurement of $\alpha_{\mathrm{IR}}$ across all observational channels would provide
strong evidence for the global CPT-symmetric architecture developed in Companions~I–IV,
while any significant mismatch would falsify the framework.

Companion~V therefore completes the observational synthesis of RDCC and establishes the
program as a unified, predictive, and empirically testable description of cosmic evolution. The
relaxation-driven loss of global coherence, quantified by $\alpha_{\mathrm{IR}}$, emerges as the single
parameter linking the earliest moments after the bounce to the largest observable scales of the
Universe.

\appendix
\section*{Appendix A — Minimal Python Implementation of the RDCC Tensor Spectrum}
\addcontentsline{toc}{section}{Appendix A — Minimal Python Implementation of the RDCC Tensor Spectrum}

This appendix provides a minimal Python implementation of the RDCC tensor spectrum,
including the characteristic log-periodic modulation arising from interference between the
two CPT-related tensor sectors across the bounce. The code is intended as a lightweight
reference implementation for reproducing the qualitative features of the RDCC stochastic
gravitational-wave background (SGWB). A full numerical pipeline, including likelihood
evaluation and parameter inference, is provided in the accompanying Companion Papers.

\begin{verbatim}
import numpy as np

# RDCC tensor spectrum with log-periodic modulation
def Omega_GW(k, alpha=0.017, eps=0.1, omega=3.2, kb=1e-3, phi=0.0):
    """
    Minimal RDCC tensor spectrum.
    Parameters:
        k     : comoving wavenumber
        alpha : relaxation parameter
        eps   : modulation amplitude
        omega : logarithmic frequency
        kb    : bounce scale
        phi   : phase shift
    """
    # Blue-tilted background spectrum
    Omega0 = k**(2*alpha)

    # Log-periodic modulation
    modulation = 1.0 + eps * np.cos(omega * np.log(k/kb) + phi)

    return Omega0 * modulation

# Example usage:
k = np.logspace(-4, 1, 500)
spectrum = Omega_GW(k)
\end{verbatim}


\end{document}
